\section{C14.1b: every state when the first testing angle is a unit}
\label{sec:r18-sphere-unit-u}

Retain the entire original C14.1b source intersection and the exact
SP6 parent from revision 17.  Assume $u$ is a bounded nonzero unit
and $w$ is infinitesimal.  The argument below proves the whole parent
for every testing state, including the previously unadmitted
$W/Z$-unbounded states.  It does not cover infinitesimal $u$.
All source gaps, $D=0$, both source sphere endpoints and the literal
coefficient boundaries remain unchanged.

\paragraph{No state chart is needed.}
All roots below are homogeneous polynomials in the original testing
state, with pure-source coefficients.  In particular the proof
uses no state-coordinate inverse, no bound on $W/Z$, and no sign
restriction.  The only testing inverse is that of the unit $u$.
Write $U=u^2$, $s=x/\sigma$, and preserve the exact SP/BS coefficients
\[
 L_h=\lambda\sigma/h,\quad \rho=\lambda x/\sigma,
 \quad d=4D/\sigma^2,\quad K_A=4\kappa h/\sigma-\lambda.
\]
The source identities give $s\in\mathrm{Inf}$,
$\sigma/\sqrt h\in\mathrm{Inf}$, $L_h\sim2$, $\nu\sim1$, and
\[
 d+\rho\sim2,\quad K_A-(2-4h)\in\mathrm{Inf},\quad
 \chi-(4-4h)\in\mathrm{Inf},\quad
 (c_G+zK_r)-2(1+U)\in\mathrm{Inf}.
 \tag{SU2}
\]
The last infinitesimal is a finite separated expression for bounded
$u$ and infinitesimal $w$.  For the only new source assertion in
SU2, the sphere identities give exactly
\[
 d=\frac{2\mathcal H}{PJ}d_0^2,\qquad
 \rho=\frac{2xH_0}{\sigma B}
       +\frac{2\mathcal H}{B}b_0^2+\lambda_0s.
 \tag{SU3}
\]
Thus $d+\rho\sim2$ follows from $d_0^2+b_0^2=1$, with no inverse
of either sphere coordinate.  All other bounds are the reviewed
SP9--SP11 and BS1 source identities.

\paragraph{Shift the troublesome angular product into the normal block.}
Set $A_0=qW$, $A=A_\theta=(U-x)A_0$, and retain $R$.
There is an exact conductor:
\[
 A(A+uR)=u(U-x)wWZ-x(U-x)A_0^2.
 \tag{SU4}
\]
Consequently the angular/radial part of BS2 is
\[
\begin{split}
 K_AA^2+\chi uAR+(c_G+zK_r)R^2
 ={}&2(A+uR)^2+2R^2-4hu(U-x)wWZ\\
 &+4hx(U-x)A_0^2+\delta_A A^2+
   \delta_\chi uAR+\delta_RR^2,
\end{split} \tag{SU5}
\]
where the three $\delta$ coefficients are the infinitesimals in
SU2.  This step keeps the signed product exactly; it is not an
inequality based on the sign of $U-x$.

\paragraph{A complete strict normal block.}
Introduce the legal separated roots
\[
 H=\sqrt h\,uZ,\qquad S=uwW-hu^2Z,\qquad N=N_x,
 \qquad Z_g=uwW=S+u\sqrt h\,H.
\]
Their exact allocations are
\[
\begin{split}
 \Gamma&=-S+\left(\frac{\sigma}{u\sqrt h}-u\sqrt h\right)H,\\
 P_h:=\sqrt h P_\beta&=(1-sh)H-\frac{s\sqrt h}{u}S,\\
 X_e&=N-\sqrt h P_h.
\end{split} \tag{SU6}
\]
Every coefficient is bounded; the factors $u^{-1}$ are pure testing
units.  The full signed packet in BS2 also satisfies exactly
\[
 \mathscr P=L_hP_h^2+
      \rho(1-x/U)Z_g^2-\rho\sigma(U-x)A_0^2.
 \tag{SU7}
\]
The apparently signed factor $1-x/U$ is simply a bounded Gram
coefficient of residue one.  No mixed inverse is used.

Combine the $wWZ$ term in SU5 with every $N,\Gamma,P_h,Z_g$ term:
\[
\begin{split}
 \mathscr N={}&[2(1+w^2)+4mw^2]N^2
     +2(2D+1)U X_e^2+d\Gamma^2-4u\nu N\Gamma\\
 &+L_hP_h^2+\rho(1-x/U)Z_g^2
       -4\sqrt h\,\frac{U-x}{u}\,Z_gH .
\end{split} \tag{SU8}
\]
This is an exact equality for the entire indicated normal block.
Its coefficients in $(N,S,H)$ are bounded separated expressions.
Using SU2 and SU6, its residue is exactly
\[
 2(1+u_0^2)N^2+4u_0NS+2S^2+2H^2.
 \tag{SU9}
\]
The $N,S$ determinant is $4$, so this residue is strictly positive.
The $H$ cross terms cancel: the $X_e$ square cancels the leading
$N\Gamma$ coupling to $H$, and the angular shift in SU5 cancels
the remaining $S H$ and $hH^2$ terms.  This is where the sphere
sum $d_0^2+b_0^2=1$ is essential.

The remaining angular/radial block is
\[
\begin{split}
 \mathscr A={}&2\{(U-x)A_0+uR\}^2+2R^2\\
 &+(4hx-\rho\sigma)(U-x)A_0^2
   +\delta_A(U-x)^2A_0^2
   +\delta_\chi u(U-x)A_0R+\delta_RR^2.
\end{split} \tag{SU10}
\]
Its residue on $(A_0,R)$ has determinant $4u_0^4>0$.

\paragraph{Every other original-parent term.}
The unchanged hard $N_y$ diagonal has residue $2(1+u_0^2)$.
All SP5 terms other than its two $\Gamma$ terms already have
infinitesimal coefficients on $(N,N_y,A_0,R)$, by SP's complete
opening-uniform coefficient bounds and $A=(U-x)A_0$.
The two $\Gamma$ terms have a factor $qw$ and become infinitesimal
crosses after the bounded allocation SU6.  Thus the complete Gram
on $(N,N_y,S,H,A_0,R)$ has strictly positive residue, the direct
sum of SU9, the $N_y$ entry and SU10.

The remaining exact nonnegative rows are retained once: both
centered cap excesses, $4\beta D_wN^2$, and
$4bpzU\Gamma^2$.  The $d\Gamma^2$ part and
$2(2D+1)(uX_e)^2$ have already been used in SU8 and are not
retained a second time.  All scalar, hard and first-normal terms
of the original SP6 parent are therefore included.

The bounded separated-Gram criterion proves tensor SOS for all real
states and both testing signs, including literal $Z=0$ and $w=0$.
Applying the real identity separately to real and imaginary parts
gives all complex states.  Neither source sphere endpoint is inverted.
The precise remaining C14.1b testing domain has $u$ infinitesimal
(with the earlier state and sign admissions still retained).

The focused checker verifies SU4, SU6--SU8 and the whole reassembled
SP6 parent, then the two strict residue determinants.  It does not
replay the unchanged source formulas or claim a whole-source closure.
